% ===================================================================
%  TABELL Nr. 8 — D'Recolte.
%  Traduction 2026 d'apres texto/io, controlee sur texto/fr, texto/de
%  et texto/nl. Consigne : voir texto/lb/10-tabell-01.tex.
%
%  LA VENDANGE DE LA SECONDE SCENE EST LE PREMIER ENDROIT DU LIVRET
%  OU CETTE COLONNE EST RICHE SANS RIEN DEVOIR A PERSONNE. Le
%  Luxembourg a ses vignes sur la Moselle, et la page se lit comme
%  celle du tcheque ou du galicien : « Wéngert » le vignoble,
%  « Wäistack » le cep, « Drauwen » les raisins, « Wënzer » le
%  vigneron, « Butt » la hotte, « Faass » le tonneau, « Douwen » les
%  douves, « Rëffer » les cercles, « Bidden » la cuve, « Moscht » le
%  mout, « Fassbënner » le tonnelier.
%
%  LE LITUANIEN AVAIT APPRIS QU'UNE LANGUE EST RICHE LA OU SON PAYS
%  L'EST — dense au bois, pauvre a la montagne. LE LUXEMBOURGEOIS LE
%  CONFIRME PAR UN AUTRE PAYS ET UN AUTRE OBJET : il n'a qu'un mot
%  pour la mer, qu'il n'a pas, et onze pour le vin, qu'il fait.
%
%  ET LA MOISSON EN DIT AUTANT : « Garwen » les gerbes, « Ähren » les
%  epis, « Halmen » les chaumes, « Seesen » les faux, « Séchel » la
%  faucille, « Wetzsteen » la pierre a aiguiser, « Dreschflegelen »
%  les fleaux, « Stréi » la paille. Ne viennent du dehors que la
%  « Méimaschinn », la « Dreschmaschinn » et le « Lokomobil » —
%  c'est-a-dire les trois machines, et rien d'autre. La partition du
%  tableau 12 lituanien, ou tout le mouvement etait nomme et tout le
%  materiel emprunte, se retrouve ici sur un champ de ble.
%
%  L'ECART DECLARE DU BLOC c2-03-1 EST RENDU COMME DANS LES QUATORZE
%  COLONNES PRECEDENTES : le fac-simile ido ecrit « gobleto 13) »
%  sans la parenthese ouvrante, mais la coquille est du compositeur
%  de l'ido et non du texte luxembourgeois, et une traduction n'est
%  pas une transcription. La parenthese est retablie.
%
%  UN HOMOGRAPHE ASSUME. « Kar » est le seigle de l'alinea 10 et
%  aussi la charrette du tableau 5 ; le luxembourgeois ecrit les deux
%  pareil, et le contexte les separe comme il le fait pour tout le
%  monde. On ne corrige pas une langue pour eviter une ambiguite
%  qu'elle assume.
% ===================================================================

%%K t08-noto-1 noto
\VUnotes{143.2mm}{%
\textit{(*) Well dat Wuert net genee ass: ass et Läinrecolte, Drauwerecolte,
Apelrecolte? Vläicht wier et gutt, „} \VUgras{mesono} \textit{“ ze benotzen, dat a
Mondo XIV, S.~242 virgeschloe gouf. Mä bis elo ass déi nei Wurzel net vun der
Akademie ugeholl a waart op d'Diskussioune no de gewéinlechen Idisteregelen.}%
}

%%K t08-apar-1 apar
\VUsaut{5.0mm}
\VUtitre{15.0pt}{60}{TABELL Nr. 8}
\VUsaut{3.0mm}
\VUcentre{13.2pt}{90}{\textit{Éischt Zeen.}}
\VUsaut{3.0mm}
\VUpk{10.2pt}{40}{D'Recolte (*).}
\VUsaut{4.0mm}
\VUpk{10.2pt}{40}{D'Biller vum Land.}
\VUsaut{3.0mm}

%%K t08-c1-01-1 p
1. --- Mam \VUgras{Summer} huet d'Bild vum Land nach eng Kéier gewiesselt. De Somen,
deen der Fuerch uvertraut gouf, ass gekeimt; e grénge Pflänzchen ass aus dem Buedem
komm an huet Ähren ugesat. D'Kar huet sech gebilt, duerno ass et reif ginn, a vun
elo un brauch een et nëmmen nach ze recolteieren.

%%K t08-c1-02-1 p
2. --- D'\VUgras{Feldblummen} sinn opgaangen. \VUgras{Kornblummen}
\textsuperscript{(1)}, \VUgras{Mohnblummen} \textsuperscript{(2)} a
\VUgras{Fangerhitt} \textsuperscript{(3)} mëschen an déi eenheetlech Faarf vun de
Weessfelder de frouen Kontrast vun hire frësche \VUgras{Bléieblieder.}

%%K t08-c1-03-1 p
3. --- D'Uebstbeem hu sech mat Blieder geschmückt; d'Uebst ass reif ginn. De klenge
\VUgras{Kiischtebam} \textsuperscript{(4)} ass voller schéine \VUgras{Kiischten}
\textsuperscript{(5)}, déi d'Kanner mat grousser Freed pléckelen an iessen.

%%K t08-c1-04-1 p
4. --- Vun elo un kann d'Vieh de ganzen Dag an der fräier Loft verbréngen.

%%K t08-c1-04-2 p
D'\VUgras{Lämmer} \textsuperscript{(6)} si grouss ginn a sinn zu schéine
\VUgras{Schof} \textsuperscript{(7)} a schéinen \VUgras{Hämmelen}
\textsuperscript{(8)} mat dichtem Wollpelz ginn. Si weiden rouheg d'\VUgras{Gras} vun
der \VUgras{Weed} \textsuperscript{(9)}, déi mat \VUgras{Margréitercher} gesprenkelt
ass.

%%K t08-c1-04-3 p
Geschwënn ginn \VUgras{dës} Béischte geschoueren, fir hir \VUgras{Woll} ze kréien,
aus där d'Duch vun eise Kleeder gemaach gëtt.

%%K t08-tit-2 sub
\VUsaut{5.0mm}
\VUpk{10.2pt}{40}{De Schäfer.}
\VUsaut{3.4mm}

%%K t08-c1-05-1 p
5. --- De \VUgras{Schäfer} \textsuperscript{(10)} schéngt net vill op si opzepassen.
Him këmmert nëmmen dat Gewidder, dat um Horizont laanschtzitt, an de
\VUgras{Schofshidder} \textsuperscript{(13)}, deen eng \VUgras{Feldfläsch}
\textsuperscript{(14)} un d'Lëpse setzt, schéngt nach méi suerglos.

%%K t08-c1-06-1 p
6. --- D'Hëtzt ass dréckend; well de \VUgras{Hond} \textsuperscript{(15)} kënnt och,
fir säin Duuscht ze läschen; hie leckt dat frëscht Waasser vum \VUgras{Baach}
\textsuperscript{(16)} a jeet en aarmt \VUgras{Fräschelchen} \textsuperscript{(17)}
op, dat sech am Ufergras verstoppt hat. Dat sprangt an dat kloert Waasser, wou
d'\VUgras{Séirosen} \textsuperscript{(18)} bléien.

%%K t08-c1-07-1 p
7. --- Eng \VUgras{Baachstëlz} \textsuperscript{(19)} bued sech nieft der
\VUgras{Quell} \textsuperscript{(20)}, déi tëscht zwee Fielsen erausspruddelt, a
verstoppt sech ënnert de \VUgras{Dornebëscher} \textsuperscript{(21)} an de
\VUgras{Hagebutzhecken} \textsuperscript{(22)}.

%%K t08-tit-3 sub
\VUsaut{5.0mm}
\VUpk{10.2pt}{40}{D'Recolte.}
\VUsaut{3.4mm}

%%K t08-c1-08-1 p
8. --- Fir d'Kanner vum Land ass d'Weessrecolte eng ganz ameséiereg Zäit. Si
\VUgras{maache frou Purzelbeem} \textsuperscript{(24)} op de \VUgras{Stréihaufen}
\textsuperscript{(23)}, si hëllefen de \VUgras{Schnitter} \textsuperscript{(26)}, a
während déi dat \VUgras{Weessfeld} \textsuperscript{(27)} mat hire \VUgras{Seesen}
\textsuperscript{(25)} méien, déi mat dem \VUgras{Wetzsteen} \textsuperscript{(30)}
gutt geschäerft sinn, sammele si de Weess a bannen en zu \VUgras{Garwen}
\textsuperscript{(31)} oder zu \VUgras{Bierdelen} \textsuperscript{(32)}.

%%K t08-c1-09-1 p
9. --- Heiansdo dierfen déi Gréisst vun hinne mat enger \VUgras{Séchel}
\textsuperscript{(12)} déi \VUgras{gëllen Halmen} \textsuperscript{(28)}
ofschneiden, déi déi schwéier \VUgras{Ähren} \textsuperscript{(29)} voller Kär droen;
an déi Kleng, déi op d'\VUgras{Herd} \textsuperscript{(33)} oppassen, déi op der
\VUgras{Weed} \textsuperscript{(34)} am Schiet vun der grousser \VUgras{Lann}
\textsuperscript{(35)} weid, fänke mat hirem \VUgras{Netz} \textsuperscript{(36)} déi
schéi \VUgras{Päiperlecken.}

%%K t08-c1-09-2 p
E puer klammen esouguer op de \VUgras{Won} \textsuperscript{(37)}, fir d'Garwen ze
schichten, déi een hinne mat enger \VUgras{Fork} \textsuperscript{(38)}
eropräicht.

%%K t08-c1-10-1 p
10. --- Op der ganzer \VUgras{Ebene} \textsuperscript{(39)}, wou d'\VUgras{Lerchen}
\textsuperscript{(41)} opfléien, herrscht Bewegung. All Leit sinn mat der Recolte
beschäftegt. Mä während déi Aarm weider déi al Geräter benotzen ---
\VUgras{Seesen, Sécheln} a \VUgras{Dreschflegelen} \textsuperscript{(40)} ---,
loossen déi räich Besëtzer hire Weess oder hir \VUgras{Kar} \textsuperscript{(43)} mat
enger \VUgras{Méimaschinn} \textsuperscript{(42)} méien. Duerno, well si d'Garwen net
an d'Scheier ze droe brauchen, maache si grouss \VUgras{Garwenhaufen}
\textsuperscript{(44)}.

%%K t08-c1-11-1 p
11. --- Da gëtt eng \VUgras{Dreschmaschinn} \textsuperscript{(47)}
erugefouert, déi e \VUgras{Lokomobil} \textsuperscript{(45)} mat engem
\VUgras{Transmissiounsrimm} \textsuperscript{(46)} dreift, an esou gëtt de Kär vum
\VUgras{Stréi} getrennt, ouni d'Feld ze verloossen, wou en gesat gouf.

%%K t08-tit-4 sub
\VUsaut{5.0mm}
\VUpk{10.2pt}{40}{D'Gewidder.}
\VUsaut{3.4mm}

%%K t08-c1-12-1 p
12. --- Dacks kommen \VUgras{Gewiddere} \textsuperscript{(53)} an der Zäit vun der
Recolte. Fir d'éischt héiert een d'Rullen vum \VUgras{Donner} vu wäitem; e
\VUgras{Westwand} \textsuperscript{(54)} fänkt un ze blosen a béit keuchend déi
décksten Beem vum \VUgras{Bësch} \textsuperscript{(49)}. Schwaarz
\VUgras{Wollécken} \textsuperscript{(56)} stiermen den Himmel; si gi vun de
blendege Zickzacke vum \VUgras{Blëtz} \textsuperscript{(55)} duerchzunn. De
\VUgras{Reen} an heiansdo den \VUgras{Hagel} fänken un ze falen a plätzen d'Recolte
nidder, déi zum Méie prett ass.

%%K t08-c1-13-1 p
13. --- Dacks setzt de Blëtz e Gebai a Brand. Esou gouf d'lescht Joer d'Daach vun
der \VUgras{Wandmillen} \textsuperscript{(50)} zu Äsche gemaach an zwee vun hire
\VUgras{Flilleken} \textsuperscript{(51)} zerbrach.

%%K t08-c1-14-1 p
14. --- No e puer Stonne gëtt et nees roueg. E glänzende \VUgras{Reebou}
\textsuperscript{(52)} erschéngt um Horizont, an d'Natur hëlt hire Liewen nees op, deen
e puer Ablécker ënnerbrach war. D'Bauere, déi sech an d'\VUgras{Hütten}
\textsuperscript{(48)} geflücht hunn, fänken nees u ze schaffen a bemeine sech mutteg,
de Schued ze flécken, deen si grad erlidden hunn.

%%K t08-tit-5 sub
\VUsaut{6.0mm}
\VUcentre{13.2pt}{90}{\textit{Zweet Zeen.}}
\VUsaut{3.6mm}

%%K t08-tit-6 sub
\VUpk{10.2pt}{40}{D'Juegd. --- D'Drauwerecolte.}
\VUsaut{1.4mm}
\VUpk{10.2pt}{40}{D'Fëscherei.}
\VUsaut{4.0mm}

%%K t08-c2-01-1 p
1. --- Ëm d'Enn vum \VUgras{August} oder d'Mëtt vum \VUgras{September,} no de
Géigenden, fänkt d'Juegd un. Kaum hunn déi éischt Sonnestrale d'\VUgras{Eidechsen}
\textsuperscript{(2)} aus hirem Lach erausgelackelt, do mécht sech de
\VUgras{Jeeër} \textsuperscript{(5)} mat senge \VUgras{Hënn} \textsuperscript{(4)} op
de Wee.

%%K t08-c2-02-1 p
2. --- Hien huet säin décke \VUgras{Juegdkostüm} \textsuperscript{(9)} aus dickem
Duch ugedoen, huet liederner \VUgras{Guêtren} ugezunn, mat deenen hien am fiichte
Gras goe kann, wou d'\VUgras{Kréiten} \textsuperscript{(1)} sech verstoppen; hien
huet seng \VUgras{Flënt} \textsuperscript{(6)} a seng \VUgras{Juegdtäsch}
\textsuperscript{(10)} geholl, an do geet hien.

%%K t08-c2-03-1 p
3. --- Den Äifer vun der Verfollegung bréngt hien an d'Mëtt vun engem Wéngert, wou
d'Drauwerecolte ganz liewweg ass. D'Kanner vum Wënzer, déi normalerweis d'Geessen um
Wee hidden, loosse si haut op gutt Gléck weiden, an, nodeems si e ale
\VUgras{Geessbock} \textsuperscript{(3)} un e Piquet gebonnen hunn, deen net verpasse
géif, seng Fräiheet zum Ausbüxen ze notzen, ginn si sech och hiren Deel um
Vergnügen. Deen een hëlt eng Drauwekëpp aus enger \VUgras{Recoltebutt}
\textsuperscript{(12)} an ameséiert sech, de \VUgras{Saaft} \textsuperscript{(14)} an
e \VUgras{Becherchen} \textsuperscript{(13)} lafen ze loossen, fir Moscht ze maachen
(ongeguerene Wäin). En anere setzt sech e \VUgras{Blieder-Kranz}
\textsuperscript{(15)} op, deen hie grad geflecht huet.

Nieft hinne kommen dreest \VUgras{Spatzen} \textsuperscript{(16)} bis virun d'Press,
fir d'Drauwen ze klauen.

%%K t08-c2-04-1 p
4. --- Während d'Kanner sech ameséieren, schaffen d'Männer. D'\VUgras{Wënzer}
\textsuperscript{(18)} droen d'Drauwen an de Keller mat \VUgras{Réckbutten}
\textsuperscript{(17)}; e \VUgras{Fassbënner} \textsuperscript{(19)} flécht
d'\VUgras{Fässer} \textsuperscript{(20)}, kontrolléiert, ob d'\VUgras{Douwen}
\textsuperscript{(22)} an d'\VUgras{Béid} \textsuperscript{(21)} an der Rei sinn, an
ersetzt déi gebrach \VUgras{Rëffer} \textsuperscript{(23)}. Hien huet och un de
\VUgras{Bidden} \textsuperscript{(25)} geschafft, an déi een d'\VUgras{Drauwen}
\textsuperscript{(26)} geheit, fir se gueren ze loossen.

%%K t08-c2-05-1 p
5. --- Wann d'Gierung fäerdeg ass, zitt ee de Wäin of a leet de Rescht op
d'\VUgras{Press} \textsuperscript{(28)}, a lues mat der \VUgras{Schrauf}
\textsuperscript{(29)} pressend dréckt ee de Saaft eraus, dee an eng
\VUgras{Kuve} \textsuperscript{(32)} leeft, an nach eng Kéier kritt ee
\VUgras{Wäin} \textsuperscript{(31)}.

%%K t08-tit-8 sub
\VUsaut{6.0mm}
\VUpk{10.2pt}{40}{Béier a Cidre.}
\VUsaut{3.4mm}

%%K t08-c2-06-1 p
6. --- Un der Mauer vum \VUgras{Keller} \textsuperscript{(27)} klëmmt en Zweig
\VUgras{Hopfen} \textsuperscript{(30)}. Do ass en nëmmen eng Zierplanz, mä a
verschiddene Länner, wou de Wäistack ganz rar ass, gëtt den Hopfen ugebaut, fir
\VUgras{Béier} ze maachen.

%%K t08-c2-07-1 p
7. --- De klengen \VUgras{Apebam} \textsuperscript{(33)} ass mat Äppel gelueden, déi,
wa si reif sinn, ganz gudde \VUgras{Cidre} ginn. An hirem \VUgras{Kefeg}
\textsuperscript{(34)} kucken de \VUgras{Kanaresvull} \textsuperscript{(35)} an de
\VUgras{Stigelitz} \textsuperscript{(36)} mat näidesche Aen op d'Spatzen, déi fräi
ronderëm dollen.

%%K t08-c2-08-1 p
8. --- Bis wäit ewech, um Hang vun den Hiwwelen, stinn a Reien
\VUgras{Wäistéck-Päl} \textsuperscript{(40)}, déi déi wonnerschéin
\VUgras{Wäistäck} \textsuperscript{(39)} halen, déi mat schwéiere \VUgras{Këpp}
gelueden sinn.

%%K t08-c2-08-2 p
D'ganzt Joer huet de \VUgras{Wënzer} \textsuperscript{(42)} geschafft, fir säi
\VUgras{Wéngert} \textsuperscript{(38)} ze këmmeren, an elo gëtt hie fir seng Méi mat
enger schéiner Recolte belount. D'\VUgras{Wënzerinnen} \textsuperscript{(41)} fëllen
d'Bidden, déi op de Won gestallt sinn, dee vu \VUgras{Mailen} \textsuperscript{(43)}
gezunn gëtt, an all Leit si frou.

%%K t08-tit-9 sub
\VUsaut{5.0mm}
\VUpk{10.2pt}{40}{D'Fëscherei.}
\VUsaut{3.4mm}

%%K t08-c2-09-1 p
9. --- Datselwecht gëllt fir d'\VUgras{Angelfëscher} \textsuperscript{(44)}, déi
vu moies un sech mat hiren \VUgras{Angelruten} \textsuperscript{(45)} um Ufer vum
Waasser néiergelooss hunn.

%%K t08-c2-09-2 p
De \VUgras{Fësch} \textsuperscript{(47)} bäisst op den \VUgras{Hoken,} soubal si
hire \VUgras{Fuedem} \textsuperscript{(46)} an d'Waasser lackelen, a kaum hunn si
Zäit, de \VUgras{Kööder} ze erneieren, do zappelt schonn en neit Affer um Enn vum
Fuedem.

%%K t08-c2-09-3 p
Trotzdeem fänkt de \VUgras{Netzfëscher} \textsuperscript{(48)}, deen aus sengem
\VUgras{Boot} \textsuperscript{(49)} d'\VUgras{Wurfnetz} an d'Mëtt vum
\VUgras{Floss} \textsuperscript{(50)} geheit, méi Fësch.

%%K t08-c2-10-1 p
10. --- Den Hierscht ass d'Saison vun de schéine Spadséiergäng: zu Fouss, zu
\VUgras{Päerd} \textsuperscript{(52)}, mam \VUgras{Vëlo} \textsuperscript{(53)}, mam
\VUgras{Auto} \textsuperscript{(54)}. D'\VUgras{Weeër} \textsuperscript{(51)} si
propper, an ee notzt déi lescht schéin Deeg, well do sammele sech schonn
d'\VUgras{Schwalwen} \textsuperscript{(56)} op den Diecher, fir mat alle Flilleken an
déi méi waarm Länner ze fléien. D'Blieder vum alen \VUgras{Nossbam}
\textsuperscript{(57)} ginn giel, d'\VUgras{Nëss} falen, an déi héich
\VUgras{Beem,} déi wéi e \VUgras{Bouquet} \textsuperscript{(58)} op der
\VUgras{Héicht} \textsuperscript{(59)} zesummestinn, verléiere geschwënn hir
Blieder.

%%K t08-c2-11-1 p
11. --- D'\VUgras{Sonn} \textsuperscript{(60)} geet méi spéit op, d'Deeg gi méi
kuerz, e zimlech schaarft Killt fänkt un, sech de \VUgras{Moien} spiere ze loossen, an
do verloosse schonn d'Flich vu \VUgras{Schnepfen} \textsuperscript{(61)} eis
ongaaschtlech Klimae.

\VUsaut{6.0mm}
\VUfilet{22mm}
